\documentclass[11pt]{article}

\usepackage[margin=1in]{geometry}
\usepackage{amsmath,amssymb}
\usepackage{booktabs}
\usepackage{hyperref}

\title{Foundations-Cleanup Synchronization Release for HAOS-IIP}
\author{Tomislav Rupi\'c \\ Independent Researcher}
\date{April 2026}

\begin{document}
\maketitle

\begin{abstract}
This paper freezes release \texttt{50.2} of HAOS-IIP as a small foundations-cleanup synchronization layer immediately after \texttt{50.1}. The purpose is narrow and explicit. Release \texttt{50.1} synchronized the public-state comparison and the broader continuum / derivation tranche, but it still listed three bounded foundational cleanup items as open: the incidence-normalization caveat (\texttt{F2}), channel completeness and the no-extra-bridge rule (\texttt{F3}), and the Hilbert-complex upgrade of the finite-dimensional sector theorem (\texttt{F4}). Release \texttt{50.2} records that those three items are now closed in bounded note form. No new numerics are introduced, no disorder-threshold outputs are frozen, and no stronger continuum claim is made. The role of this paper is simply to state what changed, what did not change, and why the next clean frontier is numerical stress-testing of the already-closed theorem ladder rather than more internal foundations hygiene.
\end{abstract}

\tableofcontents

\section{Scope}

Release \texttt{50.2} is intentionally small.

It does not add:
\begin{itemize}
\item new scalar numerics;
\item new vector numerics;
\item new disorder-threshold result bundles;
\item stronger continuum claims.
\end{itemize}

It does add:
\begin{itemize}
\item a bounded closure of \texttt{F2};
\item a bounded closure of \texttt{F3};
\item a bounded closure of \texttt{F4};
\item spine and cross-reference updates so the theorem ladder reads continuously.
\end{itemize}

\section{What Was Still Open in 50.1}

Release \texttt{50.1} was honest about three remaining foundational cleanup items:
\begin{enumerate}
\item incidence normalization caveat (\texttt{F2});
\item channel completeness and the no-extra-bridge rule (\texttt{F3});
\item Hilbert-complex upgrade of the finite-dimensional sector theorem (\texttt{F4}).
\end{enumerate}

Those were not cosmetic loose ends. They marked the points where the theorem ladder still depended on bounded caveats that had not yet been packaged as explicit standalone refinements.

\section{What 50.2 Closes}

\subsection{F2: canonical compatibility normalization}

The \texttt{T4} note had already proved the exact no-skip law
\[
d_{k+1}Q_{k+1}^+d_k=0.
\]
What remained open was whether the familiar chain identity was only a convenient coordinate choice.

Release \texttt{50.2} closes that in bounded form by showing that the middle-grade compatibility operator itself determines a unique positive square-root weight
\[
(Q_{k+1}^+)^{1/2},
\]
and that the compatibility-weighted incidence maps satisfy the chain form exactly. The remaining freedom is only unitary equivalence on the active compatibility channel.

\subsection{F3: channel completeness and no-extra-bridge}

The \texttt{T5} note had already isolated the upward and downward incompatibility channels, but it treated the exclusion of a direct bridge operator as a minimality rule.

Release \texttt{50.2} sharpens that by identifying the bounded adjacent-grade incompatibility sector with the two-channel map
\[
C_k(\omega)=\bigl(d_{k-1}^*\omega,\; d_k\omega\bigr),
\]
then showing that every positive quadratic defect on that sector is governed by one positive Gram operator on the channel space. In that representation, an explicit bridge block is only a channel-coordinate presentation of the same positive quadratic law and disappears after canonical orthogonalization.

\subsection{F4: Hilbert-complex upgrade}

The original \texttt{T6} theorem was exact in finite-dimensional inner-product spaces.

Release \texttt{50.2} upgrades that statement to Hilbert-complex language:
\begin{itemize}
\item without closed-range assumptions, the exact / harmonic / coexact split holds with closures on the image sectors;
\item with closed-range hypotheses, one recovers the exact orthogonal direct sum
\[
V_k
=
\operatorname{im} d_{k-1}
\oplus
\mathcal H^k
\oplus
\operatorname{im} d_k^*.
\]
\end{itemize}

So the finite-dimensional theorem is now embedded in a standard closed-range Hilbert-complex form without changing the bounded interpretation of the repo.

\section{Delta Table}

\begin{center}
\begin{tabular}{p{0.22\linewidth}p{0.29\linewidth}p{0.33\linewidth}}
\toprule
Item & In 50.1 & In 50.2 \\
\midrule
F2 & Open caveat after \texttt{T4} & Canonical positive-square-root normalization note \\
F3 & Open caveat after \texttt{T5} & Channel-space Gram / orthogonalization note \\
F4 & Open caveat after \texttt{T6} & Closed-range Hilbert-complex upgrade note \\
Spine status & \texttt{F2--F4} still queued & \texttt{F2--F4} marked executed in bounded form \\
Numerics & Unchanged & Unchanged \\
Vector threshold & Still open & Still open \\
\bottomrule
\end{tabular}
\end{center}

\section{What Did Not Change}

Release \texttt{50.2} does not alter the numerical authority boundary.

In particular, it does not freeze:
\begin{itemize}
\item any new disorder-threshold bundle;
\item any new transport, gap-scan, or continuum-comparison run;
\item any stronger vector-sector continuum closure claim.
\end{itemize}

So the meaning of \texttt{50.2} is not ``the vector branch is now solved.'' Its meaning is simply that the internal theorem-ladder caveats listed in \texttt{50.1} have now been packaged and closed in bounded note form.

\section{Consequence for the Next Step}

Because \texttt{F2--F4} are now closed in bounded form, the clean next frontier is no longer theorem-ladder housekeeping. The next real stress test is numerical:
\begin{itemize}
\item transported active-branch identity under refinement;
\item threshold behavior under bounded disorder;
\item universality, effective-law, and geometry / response closure on the same derived active branch.
\end{itemize}

That is exactly where the repo should spend effort next.

\section{Conclusion}

Release \texttt{50.2} is a foundations-cleanup synchronization release, not a new-physics release. It exists so that the repository record says exactly what is now true:
\begin{itemize}
\item the \texttt{T1--T8} ladder is still the core bounded derivation stack;
\item the post-\texttt{50.1} cleanup items \texttt{F2--F4} are now executed in bounded form;
\item the remaining frontier is numerical and continuum-facing, not internal theorem-ladder repair.
\end{itemize}

That is the whole point of \texttt{50.2}.

\end{document}
